\begin{theorem}
\label{XIII.2.4}
Let
% original p. 380
$f\colon X\to S$ be an $S$-scheme, $D$ a divisor on~$X$
with normal crossings relative to~$S$ \eqref{XIII.2.1},
$Y=\Supp D$, $U=X-Y$, $i\colon U\to X$ the canonical immersion. Let $F$
be a sheaf of sets (\resp of groups) on~$U$, satisfying
one of the following conditions:
\begin{enumerate}
\item[a)] $F$ is, locally for the \'etale topology \emph{on}
$X$ and on~$S$, the inverse image of a sheaf of sets (\resp of a
constructible sheaf of groups) on~$S$.
\item[b)] $F$ is locally constant constructible on~$U$ and
tamely ramified over~$X$ relative to~$S$.
\end{enumerate}
Then we have the following conclusions:
\begin{enumerate}
\item[1)]$(F,i)$ is cohomologically proper relative to~$S$
in dimension $\leq 0$ (\resp for every morphism $h\colon S'\to S$, if
$i'\colon U'\to X'$ is the inverse image of~$i$ over~$S'$, if $F'=F|U'$
and if $k=h_{(X)}$, the canonical morphism
$$
\Psi\colon k^*(\R^1_{\tame}i_*F)\to\R^1_{\tame}i_*'F'
$$
is an isomorphism).

If $F$ is a sheaf of groups prime to the residue characteristics
of~$S$ (\Ref{XIII.2.3}.b)) (tamely
ramified over~$X$ relative to~$S$), then $(F,i)$ is
cohomologically proper relative to~$S$ in dimension $\leq 1$.
\item[2)]If $F$ is a constructible sheaf of sets (\resp of groups),
$i_*F$ (\resp $\R^1_{\tame}i_*F$) is constructible.
\end{enumerate}
\end{theorem}

\subsubsection*{Proof}
For every $S$-scheme $S'$, we consider the following diagram,
all of whose squares are cartesian:
$$
\xymatrix{ U \ar[dd]_-g \ar[dr]^-i & & & U'\ar[lll] \ar[dd]_{\hbox{\raise10mm\hbox{$g'$}}}
\ar[dr]^-{i'}& \\ & X \ar[dl]_{f\mkern-8mu} & & & X'\ar[lll]_{k\hspace*{1cm}} \ar[dl]_{f'\mkern-10mu} \\ S
& & & S'\ar[lll]_-h & & }
$$
Since
% original p. 381
the question is local on~$X$ for the \'etale topology, we may
assume that $D$ is a divisor with strictly normal crossings
relative to~$S$ \eqref{XIII.2.1}; moreover, after
restricting $X$ to a neighborhood of~$Y$, we may assume $X$ smooth
over~$S$.

\Subsubsection*{Proof of~\Ref{XIII.2.4}~\textup{1)}}
\subsubsection{}
\label{XIII.2.4.1}
\emph{Case of a sheaf of sets satisfying} a). We may
assume that we have $F=g^*G$, where $G$ is a sheaf on~$S$. It
then follows from (SGA~4 XVI~3.2) that the canonical morphism
\begin{equation*}
\label{eq:XIII.2.4.1.1}
\tag{\thesubsubsection.1} f^*G\to i_*F
\end{equation*}
is an isomorphism. For every $S$-scheme $h\colon S'\to S$, we have
likewise an isomorphism $f'^*G'\to i_*'F'$; consequently the canonical
morphism
$$
\varphi\colon k^*(i_*F)\to i_*'F'
$$
is identified with the natural isomorphism
$$
k^*f^*G\simeq f'^*G'
$$

\subsubsection{}
\label{XIII.2.4.2}
\emph{Case of a sheaf of sets satisfying} b). We must
show that $\varphi$ is an isomorphism, and for this it suffices to
see that this is so at each geometric point
$\overline{x}'$ of~$X'$. Let $\overline{S}$ (\resp $\overline{X}$,
\resp $\overline{S'}$, \resp $\overline{X'}$) be the strict localization
of~$S$ (\resp~$X$, \resp~$S'$, \resp $X'$) at $\overline{x}'$ and let us put
$\overline{U}=U_{(\overline{X})}$,
$\overline{U'}=U_{(\overline{X'})}$, etc. The morphism
$\varphi_{\overline{x}}$ is identified with the canonical morphism
$$
\overline{\varphi}\colon \H^0(\overline{U},\overline{F})\to
\H^0(\overline{U'}, \overline{F'})
$$
We can find a principal covering $V$ of~$\overline{U}$, of the
type appearing in~\Ref{XIII.5.4}, such that the inverse image of
$\overline{F}$ on~$V$ is a constant sheaf with value $C$. If
$\Pi$ is the Galois group of~$V$ over~$\overline{U}$, $\Pi$
operates on~$\overline{F}|V$, and we have
\begin{equation*}
\label{eq:XIII.2.4.2.1}
\tag{\thesubsubsection.1}
\H^0(\overline{U},\overline{F}))\simeq \H^0(V,C_V)^{\Pi}
\end{equation*}
where the right-hand side denotes the set of the
elements of~$\H^0(V,C_V)$ invariant
% original p. 382
under~$\Pi$. Since $V'=V\times_{\overline{U}}\overline{U'}$ is a
principal covering of~$\overline{U'}$ with Galois group
$\Pi'\simeq \Pi$, we see that the morphism $\overline{\varphi}$
is obtained, by taking the invariants under $\Pi$, from the
canonical morphism
$$
\H^0(V,C_V)\to \H^0(V',C_{V'}).
$$
Since $V$ and $V'$ are connected \eqref{XIII.5.4}, this morphism, hence
also $\overline{\varphi}$, is an isomorphism.

Let us note that if moreover $F$ is a sheaf of groups and if $P$ is a
torsor on~$\overline{U}$ with group $\overline{F}$,
tamely ramified relative to $D$, it follows
from the preceding proof and from~\Ref{XIII.2.2} that
$(\leftexp{P}{\overline{F}},\overline{i})$ is cohomologically
proper relative to~$S$ in dimension $\leq 0$.

\subsubsection{}
\label{XIII.2.4.3}
\emph{Case of a sheaf of groups.} To show that $\Psi$ is an
isomorphism, it suffices to prove that, for every geometric
point $\overline{y}'$ of~$Y'$, the morphism
$$
\Psi_{\overline{y}'}\colon
(k^*(\R^1_{\tame}i_*F))_{\overline{y}'}\to (\R^1_{\tame}
i_*'F')_{\overline{y}'}
$$
is an isomorphism. Now, by~\Ref{XIII.2.1.2},
$\Psi_{\overline{y}'}$ is identified with the canonical morphism
$$
\overline{\Psi}\colon
\H^1_{\tame}(\overline{U},\overline{F})\to
\H^1_{\tame}(\overline{U'},\overline{F'}).
$$
Let $\widetilde{U}$ be the universal tamely
ramified covering of~$\overline{U}$ \eqref{XIII.2.1.3} and $\widetilde{F}$
the inverse image of~$\overline{F}$ on~$\widetilde{U}$. It follows
from~\Ref{XIII.5.7} in case a) and from~\Ref{XIII.5.5} in case b)
that $\H^1_{\tame}(\overline{U},\overline{F})$ is identified with the
subset $\H^1(\overline{U},\overline{F})$ formed by the classes
of~$\overline{F}$-torsors whose inverse image on~$\widetilde{U}$
is trivial. On the other hand a
classical argument (\cf IX~\Ref{IX.5}, p\ptbl\pageref{page-300})
shows that
the set of the elements of~$\H^1(\overline{U},\overline{F})$
whose inverse image on~$\widetilde{U}$ is trivial is identified
with
$\H^1(\pi_1^{\tame}(\overline{U}),\H^0(\widetilde{U},\widetilde{F})).$
We thus obtain a canonical isomorphism
\begin{equation*}
\label{eq:XIII.2.4.3.1} \tag{\thesubsubsection.1} {}
\H^1_{\tame}(\overline{U},\overline{F})\isomto
\H^1(\pi_1^{\tame}(\overline{U}),\H^0(\widetilde{U},\widetilde{F})).
\end{equation*}

Consequently the morphism $\overline{\Psi}$ is identified with the
canonical morphism
$$
\H^1(\pi_1^{\tame}(\overline{U}),\H^0(\widetilde{U},\widetilde{F}))
\to
\H^1(\pi_1^{\tame}(\overline{U'}),\H^0(\widetilde{U}',\widetilde{F}')).
$$
Let us show
% original p. 383
that this morphism is an isomorphism. The morphism
$\pi_1^{\tame}(\overline{U'})\to\pi_1^{\tame}(\overline{U})$
is an isomorphism by~\Ref{XIII.5.6}, and the same
holds for the morphism $\H^0(\widetilde{U},\widetilde{F})\to\allowbreak
\H^0(\widetilde{U}',\widetilde{F}')$. Indeed, this is obvious
in case b) since $\widetilde{F}$ is constant and $\widetilde{U}$ and
$\widetilde{U}'$ are connected. In case a), let $\overline{G}$ be a
constructible sheaf of groups on~$\overline{S}$ such that we have
$\overline{F}=\overline{g}^*\overline{G};$ since the morphisms
$\widetilde{U}\to\overline{S}$ and
$\widetilde{U}'\to\overline{S'}$ are $0$-acyclic
\eqref{XIII.5.7}, we have
$$
\H^0(\widetilde{U},\widetilde{F})\simeq
\H^0(\overline{S},\overline{G}) \simeq
\H^0(\overline{S'},\overline{G'})\simeq
\H^0(\widetilde{U}',\widetilde{F}'),
$$
which implies that $\Psi$ is an isomorphism. The last
assertion of \Ref{XIII.2.4}~1) follows from what precedes,
taking into account \Ref{XIII.2.3}~b).

\subsubsection*{Proof of~\Ref{XIII.2.4}~\textup{2)}}

The case of a constructible sheaf of sets satisfying a)
follows at once from \eqref{eq:XIII.2.4.1.1}. Let $F=g^*G$ be a
sheaf of groups satisfying a), where $G$ is a constructible
sheaf; we may assume $S$ affine; let $(S_j)_{j\in J}$ be
a finite family of reduced closed subschemes of~$S$
whose union covers $S$, such that the inverse image of~$G$ on~$S_j$ is a locally constant sheaf. Taking into account
\Ref{XIII.2.4}~1), it suffices, in order to establish that $\R^1_{\tame}i_*F$
is constructible, to see that this is so after the change
of base $S_j\to S$, for each $j\in J$. We are thus
reduced to case~b) where $F$ is locally constant.

We assume from now on that $F$ is a sheaf of sets or of
groups satisfying b). Since the question is local for the
\'etale topology on~$X$, we may assume $X$ of finite
presentation over~$S$, and, by passage to the limit, we may assume $X$
and $S$ noetherian.

Let $D=\sum_{1\leq i\leq r}\divisor f_i$, where, for each point
$x$ of~$\Supp D$, if $I(x)$ is the set of the $i$ such that
$f_i(x)=0$, the subscheme $V((f_i)_{i\in I(x)})$ is smooth over~$S$ of codimension $\card I(x)$ in~$X$. Let $\mathcal{P}$ be
the set of subsets of~$[ 1,r ]$ and, for each $I\in\mathcal{P}$,
let us put
$$
X_I=\textstyle\Big(\bigcap_{i\in I}V(f_i)\Big)\cap\Big(\bigcap_{i\not\in I}X_{f_i}\Big).
$$
Let
% original p. 384
$z$ be a point of~$X_I$. After first restricting to an
\'etale neighborhood of~$z$, we can find an open subset $W$ of~$X$
containing $z$ and a principal covering $V$ of~$U\cap W$,
tamely ramified over~$W$ relative to~$S$, of the type
considered in~\Ref{XIII.5.6.1}, such that the inverse
image of~$F$ on~$V$ is a constant sheaf with
value~$C$. Let $\pi$ be the Galois group of the covering~$V$. For
every geometric point $\overline{x}$ of~$X_I$, we then have,
by \eqref{eq:XIII.2.4.2.1},
$$
\H^0(\overline{U},\overline{F})\simeq \H^0(V,C_V)^{\pi}.
$$
It follows that $i_*F|X_I\cap W$ is locally constant (SGA~4
IX~2.13) and consequently $i_*F$ is constructible.

Let us finally show that if $F$ is a locally constant sheaf of groups,
$\R^1_{\tame}i_*F|_{X_I}$ is constructible. If
$\overline{x}$ is a geometric point of~$X$, we have obtained
in \eqref{eq:XIII.2.4.3.1} the expression
$$
\H^1_{\tame}(\overline{U},\overline{F})\isomto
\H^1(\pi^{\tame}_1(\overline{U}),\H^0(\widetilde{U},\widetilde{F})).
$$
If $p$ is the residue characteristic of~$\overline{X}$, we
have, by~\Ref{XIII.5.6},
$\pi_1^{\tame}(\overline{U})=\prod_{l\neq p }\ZZ_l(1)^{\card I}$. Let us denote by $\LL$ the set of prime numbers which
divide the order of the finite group $\H^0(\widetilde{U},\widetilde{F})$ and
let $K=\prod_{l\in\LL -\{p\}\cap\LL }\ZZ_l(1)^{\card I}$. It
follows from \cite[I \S5 ex.2]{XIII.4} that we have
$$
\H^1_{\tame}(\overline{U},\overline{F})\simeq
\H^1(K,\H^0(\widetilde{U},\widetilde{F})).
$$
Since $K$ is topologically of finite type and
$\H^0(\widetilde{U},\widetilde{F})$ finite, we first of all
deduce from this that the fibers of the sheaf $\R^1_{\tame}i_*F|_{X_I}$ are
finite. On the other hand, the set $\LL$ does not depend on the point
$\overline{x}$. For every $q\in\LL$, let $X_{I,q}$ be the closed subset of
$X_I$ with equation $q=0$ and let $X_{I'}$ be the open subset of~$X_I$
complementary to the union of the $X_{I,q}$. Then
$\R^1_{\tame}i_*F|_{X_{I,q}}$ and $\R^1_{\tame}i_*F|_{X_I'}$ are
locally constant; indeed a specialization arrow of
geometric points of~$X_{I,q}$ (\resp of~$X_{I'}$) induces an
isomorphism on the groups $K$ \eqref{XIII.5.6.1} hence also on the
sets $\H^1_{\tame}(\overline{U},\overline{F})$, and one can apply
SGA~4 IX~2.13.

\begin{corollary}
\label{XIII.2.5}
Let
% original p. 385
$f\colon X\to S$ be a morphism, $D$ a divisor on~$X$
with normal crossings relative to~$S$ \eqref{XIII.2.1},
$Y=\Supp D$, $U=X- Y$, $i\colon U\to X$ the canonical
immersion. Let $\Phi$ be a stack on~$U$ and suppose given
surjective \'etale morphisms $S_1\to S$ and $X_2\to
X\times_{S}S_1$, a stack $\Psi$ on~$S_1$ and an isomorphism
$\Phi|U_2\simeq\Psi|U_2$ (\cf \Ref{XIII.2.2}).

Then, for every morphism $h\colon S'\to S$, if $k=h_{(X)}$,
the canonical functor
$$
\varphi\colon k^*i_*\Phi\to i_*'\Phi'
$$
is fully faithful. If $\Psi$ is constructible, the canonical
functor
$$
\psi\colon k^*i_*^{\tame}\Phi\to i_*^{\prime\tame}\Phi'
$$
is an equivalence of categories.

Moreover, if the stack $\Psi$ is constructible (\resp if $\Psi$ is
$1$-constructible \eqref{XIII.0} and $S$ locally noetherian),
$i_*\Phi$ is constructible (\resp $i_*^t\Phi$ is
$1$-constructible).
\end{corollary}

Let us show that $\varphi$ is fully faithful. It suffices to see
that, for every geometric point $\overline{x}'$ of~$X'$, this
is so for the functor
$$
\overline{\varphi}\colon
\Phi(\overline{U})\to\Phi'(\overline{U'})
$$
(we have taken up again the notation of~\Ref{XIII.2.4.2}). Let $a$, $b$ be two
elements of~$\Phi(\overline{U})$, $a',\ b'$ their images under
$\overline{\varphi}$. Since the morphism
$\overline{U}\to\overline{S}$ is locally $0$-acyclic
\eqref{XIII.5.7}, we have an isomorphism
$$
\H^0(\overline{U},\overline{S\Phi})
\overset{\sim}{\to}\H^0(S,\overline{S\Psi}).
$$
Consequently $a$ and $b$ come by inverse image from elements
of~$\Psi$, and hence so does
$F=\SheafHom_{\overline{U}}(a,b)$. Since the sheaf
$F'=\SheafHom_{\overline{U'}}(a',b')$ is the inverse image on~$\overline{U'}$ of~$F$, it follows from~\Ref{XIII.2.4.1} that the
canonical morphism
$$
\H^0(\overline{U},F)\to \H^0(\overline{U'},F')
$$
is an isomorphism, which proves that $\overline{\varphi}$ is
fully faithful.

Let us show
% original p. 386
that, for every geometric point $\overline{x}'$ of~$X'$, the
functor
$$
\overline{\psi}\colon i^{\tame}_*\Phi(\overline{X'})\to
i_*^{\prime\tame}\Phi'(\overline{X'})
$$
is an equivalence. By what precedes,
$\overline{\psi}$ is fully faithful. Let us show that
$\overline{\psi}$ is essentially surjective. Let $a'$ be an
element of~$\Phi'(\overline{U'})$ tamely
ramified over~$X'$ relative to~$S'$ \eqref{XIII.2.2} and
let us show that it is the image of a tamely
ramified element of~$\Phi(\overline{U})$. It follows from~\Ref{XIII.2.4}
1) that the canonical morphism
$$
\H^0(\overline{U},\overline{S}\overline{\Phi})\to
\H^0(\overline{U'},\overline{S}\overline{\Phi}')
$$
is an isomorphism. Let $G'$ be the maximal subgerbe of~$\Phi'$
generated by $a'$; there then exists a maximal subgerbe $G$ of
$\overline{\Phi}$, inverse image of a gerbe on~$\overline{S}$, such
that we have
$$
\overline{m}^*G\simeq G',
$$
where $m$ is the morphism $\overline{U'}\to
\overline{U}$. The canonical functor
\begin{equation*}
\label{eq:XIII.2.5.*}
\tag{$*$} {\overline{k}^*\overline{i}_*^{\tame}G\to
\overline{i}_*^{\prime\tame}G'}
\end{equation*}
is an equivalence, since $G$ is identified with a gerbe of
torsors under a constructible sheaf of groups coming from
$\overline{S}$, and one can apply \Ref{XIII.2.4}~1). It
then follows from \eqref{eq:XIII.2.5.*} that there exists an
element $a$ of~$G(\overline{U})$, tamely
ramified over~$X$ relative to~$S$, whose inverse image on~$\overline{U'}$ is $a'$, which proves that $\psi$ is an
equivalence.

If $\Psi$ is constructible, so is~$i_*\Phi$; indeed
an object $x$ of~$i_*\Phi$ is, locally for the
\'etale topology of~$S$, the inverse image of an object $y$ of~$\Psi$; it
therefore follows from~\Ref{eq:XIII.2.4.1.1} that $\SheafAut (x)$ is
the inverse image of~$\SheafAut (y)$, hence is constructible. Finally, if
$\Psi$ is $1$-constructible, so is
$i^{\tame}_*\Phi$ by~\Ref{XIII.6.3} below.

\begin{corollary}
\label{XIII.2.6}
The
% original p. 387
notation is that of~\Ref{XIII.2.4}. Assume that $S$ is of
characteristic zero at every point $s$ such that we have $Y_s\neq
\emptyset.$ Then, if $\cal{ F}$ is a locally
constant constructible sheaf of groups on~$U$ (\resp a constructible stack on~$U$ which
is, locally on~$X$ and $S$, the inverse image of a constructible stack
on~$S$), $(\cal{ F},i)$ is cohomologically proper relative
to~$S$ in dimension $\leq 1.$
\end{corollary}

Since every constructible sheaf of sets on~$U$ is
tamely ramified over~$X$ relative to~$S,$ the
corollary follows from~\Ref{XIII.2.4} (\resp \Ref{XIII.2.5}).

\begin{corollary}
\label{XIII.2.7}
The notation is that of~\Ref{XIII.2.4}, but we are given in addition
an $S$-scheme $T$, a proper morphism $p\colon X\to T,$
and we assume $X$ and $T$ of finite presentation over~$S$; let
$q=pi$. Let $\cal{ F}$ be a constructible sheaf of sets on~$U$
satisfying one of the conditions \textup{a)} or \textup{b)} of~\Ref{XIII.2.4}
(\resp a sheaf of groups satisfying one of the conditions
\textup{a)}, \textup{b)} of~\Ref{XIII.2.4}, \resp a stack on~$U$ which is, locally on~$X$ and $S$, the inverse image of a constructible stack $G$ on~$S$). Then
we have the following conclusions:
\begin{enumerate}
\item[1)] $(\cal{ F},q)$ is cohomologically proper relative to~$S$ in dimension $\leq 0$ (\resp for every morphism $h\colon
S^{'}\to S$, if $m=h_{(T)}$, the canonical morphism:
$$
\Theta\colon m^*(\R^1_{\tame}q_*\cal{ F})\to
\R^1_{\tame}q_*^{'}\cal{ F}^{'}
$$
is an isomorphism, \resp for every morphism $S^{'}\to S$,
the canonical morphism
$$
\xi \colon m^*(q_*^{\tame}\cal{ F})\to
{q'}_*^{{\tame}}\cal{F}^{'}
$$
is an equivalence).
\item[2)] the sheaf $q_*\cal{ F}$ (\resp the sheaf
$\R_{\tame}^1q_*\cal{ F}$, \resp the stack $q_*^{\tame}\cal{ F}$)
is constructible. In the last case, if we assume $S$ locally
noetherian and $G$ $1$-constructible, the same is true
of $q_*^{\tame}\cal{ F}.$
\end{enumerate}
\end{corollary}

The
% original p. 388
first part follows at once from~\Ref{XIII.2.4},
\Ref{XIII.2.5} and from the proof
of~\Ref{XIII.1.8}. Let us prove~2). If $\cal{ F}$ is a constructible sheaf
of sets on~$U$ satisfying \Ref{XIII.2.4}~a)
or \Ref{XIII.2.4}~b), it follows from \Ref{XIII.2.4}~2) that
$i_*\cal{ F}$ is constructible; hence so is
$q_*\cal{ F}=p_*(i_*\cal{ F})$ (SGA~4 XIV~1.1).

Let $\cal{F}$ be a constructible sheaf of groups on~$U$
satisfying \Ref{XIII.2.4}~a) or \Ref{XIII.2.4}~b) and let us prove
that $\R_{\tame}^1q_*\cal{ F}$ is constructible. By passage to
the limit (EGA~IV 8.10.5 and 17.7.8) and using 1), we may
assume $S$ noetherian. Let then $\Phi$ be the stack on~$X$ whose
fiber at every scheme $X^{'}$ \'etale over~$X$ consists
of the torsors on~$U^{'}=U\times_{X}X^{'}$, with group $\cal{F}|U$, which
are tamely ramified over~$X$ relative to~$S$; we
therefore have
$$
S(i_*^{\tame}\Phi)\simeq \R_{\tame}^1i_*\cal{ F}
$$
and this sheaf is constructible by \Ref{XIII.2.4}~2).
It therefore follows from~\Ref{XIII.6.3} below that $S(p_*\Phi)$
is constructible, \ie that $\R^1_{\tame}q_*\cal{ F}$ is
constructible.

Finally, if $\cal{ F}$ is a stack on~$U$ which is, locally on~$X$ and
$S$, the inverse image of a constructible stack on~$S,$
$i_*^{\tame}\cal{ F}$ is constructible, and hence so
is~$q_*^{\tame}\cal{ F}=p_*(i_*^{\tame}\cal{ F})$. If
moreover $S$ is locally noetherian and $SG$ constructible,
$S(i_*^{\tame}\cal{ F})$ is constructible by 6.3; hence
so is~$S(q_*i_*^{\tame}\cal{ F})$, \ie
$S(q_*^{\tame}\cal{ F})$~\Ref{XIII.6.2}.

\begin{corollary}
\label{XIII.2.8}
Let
$$
\xymatrix{ U \ar[r]^{i} \ar[d]_{g} & X \ar[ld]^{f} \\ S }
$$
be a commutative diagram of schemes, in which $U$ is the open set
complementary
% original p. 389
in $X$ to a divisor with normal crossings relative to~$S$, $f$ a \emph {proper} morphism of finite presentation. Let
$\LL$ be a set of prime numbers. We assume~$g$ locally
$0$-acyclic (\resp locally $1$-aspherical for $\LL$). Then, if
$\cal{ F}$ is a sheaf of sets on~$U$ (\resp a sheaf of
$\LL$-groups) on~$U,$ locally constant constructible,
tamely ramified over~$X$ relative to~$S$,
$f_*\cal{ F}$ (\resp $\R^1_{\tame}f_*\cal{ F}$) is locally
constant constructible and $(\cal{ F},f)$ is cohomologically proper
relative to~$S$ in dimension $\leq 0$ (\resp the formation of
$\R^1_{\tame}f_*\cal{ F}$ commutes with every change of base
$S^{'}\to S$). In the case without ``resp.'', if $\cal{ F}$ is a
sheaf of groups, for every specialization $\overline{s}_1
\to \overline{s}_2$ of geometric points of~$S$, the
specialization morphism
$$
(\R^1_{\tame}f_*\cal{ F})_{\overline{s}_2}\to
(\R^1_{\tame}f_*\cal{ F})_{\overline{s}_1}
$$
is injective.
\end{corollary}
The corollary follows at once from~\Ref{XIII.2.6} and
from~\Ref{XIII.1.14} (\resp from the analogue of~\Ref{XIII.1.14} for
$\R^1_{\tame}f_*\cal{ F}$, which is proved as in \loccit).

\begin{corollary} \label{XIII.2.9}
Let
$$
\xymatrix{ U \ar[r]^{i} \ar[d]_{g} & X \ar[ld]^{f} \\ S }
$$
be a commutative diagram of schemes, in which $U$ is the open set
complementary in~$X$ to a divisor with normal crossings
relative to~$S$, $f$~a proper smooth morphism of
finite presentation. Let~$\LL $ be the set of prime numbers
distinct from the residue characteristics of~$S$. Let~$\cal{F}$ be a locally constant constructible sheaf of~$\LL $-groups
on~$U$, tamely ramified over~$X$ relative
to~$S$. Then $\R^1 f_*\cal{ F}$ is locally constant
constructible and $(\cal{ F},f)$ is cohomologically proper
relative
% original p. 390
to~$S$ in dimension $\leq 1$.
\end{corollary}
The corollary follows from~\Ref{XIII.2.8} and from the fact that we have
$\R_{\tame}^1f_*\cal{ F}=\R^1f_*\cal{ F}$ (\Ref{XIII.2.3}~b)).
